\title{Large Language Models Can Automatically Engineer Features for Few-Shot Tabular Learning}

\begin{abstract}
Large Language Models (LLMs) have demonstrated extraordinary capabilities in addressing complex and novel reasoning challenges, positioning them as valuable tools for tabular learning—a key component in many real-world contexts. This work introduces \textsf{FeatLLM}, a new in-context learning paradigm in which LLMs act as feature constructors, transforming input data into a representation highly amenable to tabular prediction tasks. The resulting features inform a basic downstream model, such as linear regression, thereby enabling effective few-shot learning with strong predictive outcomes. Importantly, at the time of inference, the \textsf{FeatLLM} framework utilizes only the discovered features in conjunction with this simple predictive model, dispensing with reliance on the LLM itself. In contrast to prevailing LLM-driven approaches, \textsf{FeatLLM} obviates the need to make LLM queries for each inference instance. Additionally, it requires only API access to the LLM and is unaffected by constraints on prompt size. Comprehensive experiments spanning diverse tabular datasets reveal that \textsf{FeatLLM} crafts high-caliber features, achieving a consistent average improvement of 10\% over methods such as TabLLM and STUNT.
\end{abstract}

\section{Introduction}
Large Language Models (LLMs) have recently garnered significant attention for their capacity to address novel tasks without the necessity for specialized fine-tuning~\cite{creswell2022selection,imani2023mathprompter,taylor2022galactica}. Benefiting from training on vast repositories of text, LLMs encode substantial prior knowledge that can effectively replace human expertise in a variety of fields~\cite{Genomes2021,singhal2023large,wilcox2003role}, supporting automation in areas ranging from conversational analysis~\cite{finch2023leveraging} to automated program repair~\cite{fan2023automated}. Notably, by providing an explicit task statement and a limited number of examples within their prompts, LLMs are able to grasp contextual cues and emphasize salient features and training samples~\cite{brown2020language,zhao2023survey}. This ability demonstrates the promise of LLMs as sophisticated feature engineers capable of discerning essential aspects of data. 

The exploitation of LLMs’ extensive knowledge and logical reasoning has also expanded into tabular learning scenarios. For example, recognizing that aging is a risk factor for certain diseases may enhance predictive modeling in medical tasks. Recent strategies in this space encode both feature and task specifications in natural language, linearize tabular records, and input these sequences to the LLM~\cite{dinh2022lift,wang2023anypredict}. Subsequently, approaches may employ in-context learning~\cite{nam2023semi} or pursue lightweight fine-tuning mechanisms~\cite{dinh2022lift,hegselmann2023tabllm}. These efforts reveal that leveraging LLMs' embedded knowledge yields superior outcomes compared to traditional tabular algorithms, particularly when labeled data is scarce~\cite{hegselmann2023tabllm}.

Contemporary methods that leverage LLMs for tabular data learning face several drawbacks. For full end-to-end prediction, every individual data point necessitates at least one LLM inference, which incurs significant computational overhead. Additionally, achieving optimal predictive accuracy with these methods frequently demands fine-tuning of the LLM, making them impractical when training infrastructure or tuning APIs are inaccessible—a critical limitation as leading LLMs nowadays often restrict usage to inference APIs only~\cite{anil2023palm,openai2023gpt4}. Further, the majority of existing strategies struggle with prompt length limitations~\cite{liu2023lost}; when dealing with tabular datasets featuring large numbers of columns, prompts quickly grow in size, often leading to diminished performance or outright infeasibility. Collectively, these issues complicate practical deployment of LLM-driven tabular learning solutions.

Our methodology rethinks the conventional application of LLMs, moving beyond their direct use in end-to-end prediction. Instead, we utilize them for the task of feature engineering, tapping into their expansive prior knowledge for more effective data utilization. We propose an innovative in-context learning strategy termed \textsf{FeatLLM}, designed to reveal and leverage the “criteria” underlying the LLM’s predictive reasoning. For example, in a disease prediction scenario, the LLM can be prompted to deduce and articulate explicit rules indicating which combinations of feature values signal the disease’s presence. These derived rules inform the construction of new, potentially more informative features, serving as substitutes for raw attributes in downstream tasks. This framework repositions LLMs as engines of feature creation, leading to simplified downstream modeling—once these features are extracted, there is no further need for complex machine learning architectures at inference, thereby reducing inference time. Through in-context learning, features are extracted simply by augmenting prompts with illustrative demonstrations, negating the necessity for model retraining. Moreover, by incorporating ensemble-inspired techniques such as feature bagging~\cite{breiman2001random}, our approach addresses issues arising from excessive prompt length.

\textsf{FeatLLM} decomposes the prompt-based feature engineering process into two principal phases: (1) grasping the underlying problem and inferring how features correlate with the target variables; and (2) leveraging the insights and few-shot examples from the initial step to construct critical rules for class separation. This modular reasoning guides LLMs to exclude irrelevant features during the rule formulation stage. The inferred rules are operationalized by translating them into executable code, subsequently mapping each data instance to a binary indicator reflecting its compliance with each rule. A parsimonious predictive model is then trained on these binary features to estimate class membership probabilities. To further bolster robustness, an ensemble strategy is employed, involving repeated feature generation paired with feature bagging. Tuning the selection of features and data samples within the bagging framework helps to address issues related to prompt length. As a result, \textsf{FeatLLM} avoids the burdens of retraining and imposes low computational overhead during inference, making LLM-powered feature engineering viable for a broad range of practical use cases.

Our empirical study benchmarks \textsf{FeatLLM} across 13 tabular datasets under few-shot learning scenarios, demonstrating its effectiveness and resilience. Results indicate that LLMs in our framework adeptly blend both pre-existing domain knowledge and patterns discovered from the data to produce informative features. \textsf{FeatLLM} consistently outperforms state-of-the-art few-shot learning approaches in multiple contexts. The source code is available at the anonymous repository~\texttt{https://github.com/Sungwon-Han/FeatLLM}.

\section{Related Work}

\subsection{Few-Shot Learning with Tabular Data} Developments in few-shot learning have made it possible to obtain representations that generalize well, even when very few labeled examples are available~\cite{chen2018closer}. Originally concentrated on images, few-shot learning methodologies have exhibited considerable adaptability across various modalities, now encompassing domains such as text, audio, and more recently, tabular datasets~\cite{majumder2022few,nam2023stunt,schick2021exploiting}. Nevertheless, the challenge of sparse labeled data increases susceptibility to misleading or coincidental patterns~\cite{bartlett2021deep}. In response, contemporary research has integrated supplementary information sources or incorporated domain-specific inductive biases to guide the learning process. Examples include leveraging abundant unlabeled data for semi-supervised training via data augmentation~\cite{ucar2021subtab,yoon2020vime}, or optimizing model initialization through domain-agnostic pre-training methods that require no supervision~\cite{somepalli2021saint}. Pre-training objectives in these paradigms might involve techniques such as masking or intentionally corrupting data features followed by reconstruction~\cite{arik2021tabnet,majmundar2022met}, or utilizing augmented samples for objectives based on contrastive learning~\cite{bahri2022scarf,chen2020simple}. However, because of the inherent heterogeneity in tabular datasets, designing robust, universally effective augmentation approaches remains challenging, and these strategies tend to yield limited gains for few-shot scenarios~\cite{nam2023stunt}.

Recently, approaches have surfaced that advocate pre-training on a large, heterogeneous set of tabular datasets before deploying knowledge through transfer learning to the actual task~\cite{zhang2023generative,levin2022transfer,zhu2023xtab}. For example, TransTab~\cite{wang2022transtab} trains transformer models to understand the semantics of columns, utilizing both the column names and cell values. This extracted knowledge supports both zero-shot and few-shot transfer on previously unseen tabular tasks. TabPFN~\cite{hollmann2022tabpfn}, in contrast, creates synthetic tabular data by sampling from diverse distributions to pre-train a Bayesian neural network. At inference, the model processes both the available training set and queries without further gradient updates. The accumulation of prior knowledge from such a diverse dataset corpus enables these models to outperform classic tree-based methods and delivers notable improvements in few-shot learning.

\subsection{Language-Interfaced Tabular Learning} Another line of research investigates the potential of large language models (LLMs)~\cite{anil2023palm,openai2023gpt4} as sources of auxiliary knowledge. These LLMs, trained on vast and multifaceted textual data, have demonstrated remarkable flexibility, allowing them to tackle new tasks with minimal or no further training~\cite{brown2020language}. Capitalizing on their encoded prior knowledge, several recent efforts have focused on repurposing LLMs for tabular learning tasks by converting table-structured data into text formats suitable for textual input prompts~\cite{dinh2022lift,hegselmann2023tabllm,wang2023anypredict}. Through explicit prompt engineering—which can include the tabular data itself, alongside comprehensive feature and task descriptions—LLMs are encouraged to learn and apply correspondences between features and predictive tasks. Some work has pursued approaches such as parameter-efficient fine-tuning (for instance, using LoRA~\cite{hu2021lora} or IA3~\cite{liu2022few}), while others focus on in-context learning, where representative few-shot examples are presented as prompt demonstrations, eliminating the need for explicit gradient-based training~\cite{dinh2022lift,hegselmann2023tabllm,nam2023semi,slack2023tablet}.

While prior strategies have been successful in advancing few-shot tabular learning, their dependence on routing all inference through the LLM introduces practical obstacles for industrial deployment. This paper seeks to move away from the typical black-box paradigm of relying on LLMs to predict individual samples in an end-to-end manner. Rather than processing every prediction this way, the approach here is to employ the LLM to deduce, for each class, the corresponding decision rules or defining conditions. \textsf{FeatLLM} transforms these identified rules into program code that constructs novel features, facilitating inference without recurrent LLM calls. This leverages the LLM's in-context learning capacities to derive useful rules from both few-shot examples and relevant prior knowledge.

\section{Methods}
The core purpose of \textsf{FeatLLM} is to derive the defining ``rules"—that is, the specific conditions linked to each possible answer class—from the small number of examples it receives as input, as visualized in Figure~\ref{fig:main_model}. The LLM extracts these rules, which are subsequently parsed into binary features for any given sample, representing whether each condition is met or not. These binary indicators collectively serve as inputs to a dot product layer equipped with non-negative, trainable coefficients, yielding a score that reflects class likelihood. This training procedure allows the system to calibrate the relative contribution of each feature toward class assignment (Section~\ref{Sec:training}). To enhance robustness and generalization, this procedure is performed repeatedly using bagging, and results are integrated through an ensemble. The subsequent sections provide formal definitions of the problem as well as detailed descriptions of every methodological stage.

\vspace*{-0.12in}
\paragraph{Problem formulation.} Consider a tabular dataset composed of $N$ annotated instances, $\mathcal{D} = \{(\mathbf{x}^i, \mathbf{y}^i)\}_{i=1}^{N}$. Each instance in $\mathcal{D}$ comprises $d$ input attributes (thus, $\mathbf{x}^i$ is represented as a $d$-dimensional vector), where the features have natural language identifiers, $F = \{f_j\}_{j=1}^{d}$, and each feature has an associated definition provided elsewhere. In the context of classification, the label $\mathbf{y}^i$ refers to one class among a pre-specified collection of categories. For $k$-shot learning settings, the model is trained using only $k$ ($<N$) randomly selected labeled samples.

\subsection{Prompt Design for Extracting Rules}
\label{Sec:prompt_design}
To ensure that the LLM generates rules through a more robust reasoning process, the approach is structured to reflect the systematic methodology an expert would employ for tabular data analysis. The constructed prompt is divided into three primary segments, as illustrated in Fig.~\ref{fig:prompt_for_rule} and detailed in Appendix~\ref{sec:example_rule}:

\vspace*{-0.12in}
\paragraph{Basic information description.} This component supplies all requisite background for solving the task (illustrated by the orange text in Fig.~\ref{fig:prompt_for_rule}). It articulates the challenge with a corresponding task description, includes label details, provides explanations of all features, and incorporates example demonstrations. The description of the task is posed in interrogative form (for instance: ``Does this patient’s myocardial infarction complications data indicate chronic heart failure? Yes or no?''; see additional task prompts in Appendix~\ref{sec:task_desc}). Feature descriptions specify the variable type and, if applicable, further elucidate their meaning. For categorical features, representative examples of the allowable classes may be provided. Demonstrative samples from the training set are rendered as text along with their true labels, with serialization formatted according to earlier methods~\cite{dinh2022lift,hegselmann2023tabllm}:

\begin{align}
&\text{Serialize}(\mathbf{x}^i,\mathbf{y}^i, F) = \nonumber \\ 
&``\ f_1\ \text{is}\ \mathbf{x}^i_1.\ ... \ f_d\  \text{is}\  \mathbf{x}^i_d.\ \ \text{Answer:}\  \mathbf{y}^i\ ”
\end{align}

\vspace*{-0.12in}
\paragraph{Reasoning instruction.} Our objective is to improve the LLM's ability to reason by equipping it with explicit reasoning instructions (refer to blue-highlighted text in Fig.~\ref{fig:prompt_for_rule}). The instruction process starts with a preamble modeled after the chain-of-thought methodology~\cite{wei2022chain}. The LLM's reasoning is subsequently structured into two distinct phases. Initially, the model is prompted to independently deduce causal associations or overarching patterns between input features and the overarching task description, intentionally avoiding reference to any illustrative demonstrations; this encourages reliance on domain knowledge or widely-held reasoning, thereby enabling the LLM to mine its internalized prior knowledge regarding the problem at hand. Following this, in the second phase, the LLM integrates both the illustrative demonstrations and insights from the prior step to formulate class-wise inference rules. By segmenting the reasoning procedure in this manner, the model is steered away from spurious correlations in irrelevant feature columns and prompted to emphasize salient attributes. To strike a balance between under- and over-parameterization—avoiding too few or excessively many rules and the resulting risk of underfitting or unwieldy prompts—a constant quota of rules (i.e., 10)\footnote{Discussion of rule count is presented in Section~\ref{sec:analysis}.} is imposed for each class. Moreover, during the rule formulation phase, the LLM is directed to reference the specified types for each feature as delineated in the initial description, ensuring alignment between rule logic and feature type.

\vspace*{-0.12in}
\paragraph{Response instruction.} To ensure that the inferred rules are straightforward to interpret and utilize, we offer the LLM explicit guidance on structuring its output via response instructions (see yellow-highlighted sections in Fig.~\ref{fig:prompt_for_rule}). These prompt-based instructions delineate the expected output format for each answer class and clarify the syntactic structure that each rule should embody. This structured guidance supports the LLM in selecting output formats that are coherent with the underlying feature types. In scenarios where no explicit guidance is given, the LLM possesses the capability to compose complex rules by linking several statements using logical operators (such as AND and OR). An exemplar of the resulting output is provided in Appendix~\ref{sec:outcome-rule}.

\subsection{Inferring Class Likelihood via Rules}
\label{Sec:training}

\paragraph{Parsing rules for feature generation.} The set of rules established in the earlier phase is utilized to derive new binary features. Each of these features is associated with a particular class and serves as an indicator of whether the sample adheres to the corresponding class-specific rules (with values of '0' or '1'). These derived features can then serve to estimate a sample’s membership likelihood for each class. Given that the rules created by the LLM are expressed in natural language, it is necessary to devise a parsing strategy for automating the generation of features. Although formatting constraints are imposed during rule creation to facilitate parsing, the LLM may introduce variable syntactic forms or noise~\cite{kovavc2023large}. As an illustration, expressions like ``$>$” and ``$<$” may be replaced with ``greater than” or ``less than,” respectively, contributing to the challenge of consistent parsing.

In order to address the issues of parsing non-standard or noisy textual representations, we refrain from designing intricate parser implementations and instead exploit the LLM’s capacity for semantic understanding. The LLM is well suited to interpret the intended meaning regardless of syntactic variation and is competent in generating simple code upon receiving appropriate instructions—benefiting from its code-training dataset. To take advantage of this, we construct prompts that specify the function name, expected input and output, and the rules to be encoded, then supply these prompts to the LLM. Figures~\ref{fig:prompt_for_parsing} and~\ref{sec:example_parsing}-\ref{sec:example_parse_outcome} in the Appendix provide examples of such prompts along with their respective outputs. The code returned by the LLM is executed via Python’s \texttt{exec()} function to carry out the data transformation.

\vspace*{-0.12in}
\paragraph{Inferring class likelihood.} With the binary features produced from class-specific rules, a straightforward approach to quantifying a sample's likelihood of belonging to each class is to tally the number of class rules that the sample fulfills (that is, to sum the binary indicators for each class). Nevertheless, the informativeness of rules or features is often heterogeneous, prompting the need to estimate their relevance from training data. We pursue this by utilizing a standard machine learning (ML) model, which is applied to the binary features corresponding to each class. As an example, a linear model without an intercept term could be adopted. Denote by $\mathbf{z}_k^i \in \{0, 1\}^R$ the binary feature vector for class $k$ generated from instance $\mathbf{x}^i$, with $R$ rules for each class and $c$ being the number of classes. The probability associated with each class for instance $i$, denoted as $\mathbf{p}^i$, is then calculated as follows:

\begin{align}
\text{logit}_k^i &= \max(\mathbf{w}_k, 0) \cdot \mathbf{z}_k^i, \nonumber \\
\mathbf{p}^i &= \text{Softmax}({[\text{logit}_{1}^i, ..., \text{logit}_{c}^i]}). \label{eq:infer_prob}
\end{align}

In this context, $\mathbf{w}_k$ denotes the tunable parameter associated with each feature in the linear classifier, quantifying feature relevance. By constraining $\mathbf{w}_k$ to the positive domain, we guarantee that features extracted by the LLM do not diminish a class’s predicted probability during training. Subsequently, the output logits are transformed into class probabilities via the softmax function. The weight parameters $\mathbf{w}_k$ are optimized by minimizing cross-entropy loss over a limited set of training samples.

\vspace*{-0.12in}
\paragraph{Ensembling with bagging.} To enhance prediction accuracy, the entire modeling procedure is iteratively repeated, resulting in an ensemble of inference models. Let $\{\mathbf{w}[t]\}_{t=1}^T$ represent the series of optimized weights from $T$ independent runs. The aggregated class probability for an example is then obtained by averaging over the probabilities $\mathbf{p}^i$ calculated with each $\mathbf{w}[t]$, as formalized in Eq.~\ref{eq:infer_prob}.

Randomness in rule extraction for the ensemble is fostered by configuring the LLM with a relatively high temperature parameter. Additionally, few-shot example orders are permuted for each query, based on findings that the order significantly influences LLM responses~\cite{lu2022fantastically}\footnote{Further analysis regarding the variety of generated rules is provided in Appendix~\ref{sec:diversity}}. To exploit this diversity for more robust ensemble predictions, we incorporate bagging, i.e., sampling a subset of features or training instances per trial—a strategy that becomes more crucial as the number of features or training points increases. The ensemble-based framework offers two primary strengths. First, errors arising from flawed rules in individual ensemble members can be compensated by correctly generated rules from others, thus bolstering the reliability of the final outcome; this parallels the self-consistency property of LLMs~\cite{wang2022self}, where rules consistently emerging across ensemble members tend to be more trustworthy. Second, the ensembling approach tackles limitations imposed by LLM prompt size: for datasets whose tabular representation is too large for a single prompt, bagging serves as a dimensionality reduction mechanism, ensuring consistent model performance. Consequently, even if any individual ensemble member lacks full coverage over all features or examples, the collective—composed of diverse weak learners across iterations—effectively mitigates coverage deficiencies present in single-trial rule sets.

\section{Experiments}

We assess the effectiveness of \textsf{FeatLLM} on a diverse set of tabular datasets, performing a series of analyses to elucidate the operational mechanisms of our framework. Due to space limitations, supplementary experiments—including those investigating different LLM configurations, qualitative analyses, and other aspects—are provided in the Appendix.

\subsection{Performance Evaluation}

\paragraph{Datasets.}
A total of 13 datasets, covering both binary and multi-class classification tasks, are employed in our study: (1) Adult~\cite{asuncion2007uci}, which involves determining whether an individual earns more than \$50,000 per year; (2) Bank~\cite{moro2014data}, aimed at forecasting whether a bank customer will opt into a term deposit; (3) Blood~\cite{yeh2009knowledge}, focused on predicting the likelihood of repeat blood donations; (4) Car~\cite{kadra2021well}, which categorizes car quality; (5) Communities~\cite{misc_communities_and_crime_183}, targeting prediction of regional crime rates; (6) Credit-g~\cite{kadra2021well}, identifying individuals as either good or bad credit risks; (7) Diabetes\footnote{\texttt{kaggle.com/datasets/uciml/pima-indians-diabetes-database}}, for forecasting diabetes status; (8) Heart\footnote{\texttt{kaggle.com/datasets/fedesoriano/heart-failure-prediction}}, centered on predicting the onset of coronary artery disease; and (9) Myocardial~\cite{misc_myocardial_infarction_complications_579}, for identifying cases of chronic heart failure.

Additionally, we incorporate more recent, and synthetic, datasets to ensure minimal overlap with those used during LLM pre-training: (10) Cultivars\footnote{\texttt{archive.ics.uci.edu/dataset/913}}, for predicting soybean cultivar grain yield; (11) NHANES\footnote{\texttt{archive.ics.uci.edu/dataset/887}}, which classifies a person’s age group based on their records; (12) Sequence-type, which assigns sequences to one of four categories—arithmetic, geometric, Fibonacci, or Collatz; and (13) Solution-mix, where the task is to determine if the concentration percentage in a mixture from four solutions surpasses 0.5. The first two were contributed to the UCI Machine Learning Repository after September 2023 and, as such, were not included in the pre-training phase of the LLM API (gpt-3.5-turbo-0613) that underpins our approach. The remaining two are synthetic and, consequently, cannot be present in the memory of the LLM API, allowing for unbiased evaluation.

The chosen datasets differ with respect to scale and complexity, as shown in Table~\ref{Tab:basic-desc}. Each dataset offers explicit naming and description for every feature. Notably, datasets such as ‘Communities’ and ‘Myocardial’ present over 100 attributes each, presenting additional difficulties given the limited context window available to LLMs.

\vspace*{-0.12in}
\paragraph{Baselines.}
Our experiments benchmark against nine existing methods. The initial group consists of three classical tabular modeling approaches: (1) Logistic regression (LogReg), (2) XGBoost~\cite{chen2016xgboost}, and (3) RandomForest~\cite{ho1995random}. We then examine two algorithms that can utilize unlabeled datasets: (4) SCARF~\cite{bahri2022scarf}, and (5) STUNT~\cite{nam2023stunt}. It is important to recognize that assembling large collections of unlabeled data may prove impractical in some real-world settings. Another contemporary baseline is (6) TabPFN~\cite{hollmann2022tabpfn}, which pretrains models using synthetically generated datasets of varied distributions. The final three methods leverage LLMs: (7) In-context learning (In-context)~\cite{wei2022emergent}, (8) TABLET~\cite{slack2023tablet}, and (9) TabLLM~\cite{hegselmann2023tabllm}. The In-context learning method injects a handful of training examples into the prompt at inference time with no model tuning. TABLET extends this by supplementing prompts with auxiliary information—such as rule-based hints and prototype-based features—via an external classifier to improve prediction. TabLLM adapts the LLM using parameter-efficient strategies, notably IA3~\cite{liu2022few}, for efficient few-shot learning.

\vspace*{-0.12in}
\paragraph{Implementation details.}
\textsf{FeatLLM} is constructed with GPT-3.5 serving as the foundational LLM, although the framework itself remains flexible with respect to the particular LLM employed (see Appendix~\ref{sec:backbones} for empirical results using alternative backbones). The temperature parameter during LLM inference is fixed at 0.5, while the top-p value utilizes the API default of 1. Both the ensemble size and the number of extraction rules are set at 20 and 10, respectively. Additional analysis on hyper-parameter choices can be found in Figure~\ref{fig:hyperparameter2} as well as Appendix~\ref{sec:hyper-parameter}.

For the linear model component, the Adam optimizer is applied with a learning rate of 0.01, and training proceeds for 200 epochs. We use $k$-fold cross-validation to determine the optimal training epoch. When re-implementing baseline methods, GPT-3.5 is adopted for in-context learning tasks, whereas TabLLM leverages T0~\cite{sanh2022multitask} consistent with its original configuration. It should be noted that TabLLM is confined to LLMs offering openly available checkpoints; as such, GPT-3.5 is not compatible within this baseline. In the case of traditional machine learning algorithms, namely LogReg, XGBoost, and RandomForest, parameter selection is handled by means of Grid-search in combination with $k$-fold cross-validation. The parameter $k$ is chosen as either 2 or 4 to ensure that all class labels are represented in the training split. Further implementation specifics are detailed in Appendix~\ref{sec:baselines}.

\vspace*{-0.12in}
\paragraph{Results.}
Tables~\ref{tab:main1} and \ref{tab:main2} display comparative results over various datasets. For each setting, experiments were conducted three times using different random splits of the training and test sets. We report both the mean and standard deviation of the resulting AUC scores. \textsf{FeatLLM} demonstrates robust performance, frequently ranking first or second among the evaluated baseline methods. Remarkably, with only 4-shot inputs, our approach achieves results on par with classic tabular methods trained with 32 shots (refer to Fig.~\ref{fig:compares}). Although STUNT and TabPFN provide considerable gains on selected datasets, their aggregate advantage over traditional machine learning models is modest. Furthermore, Large Language Model-based systems such as in-context learning, TABLET, or TabLLM were unable to generate predictions for tabular data characterized by a high feature count. In contrast, by employing both bagging and ensemble strategies, our framework remains effective even as the feature dimension grows, consistently achieving commendable results. It should be noted that the principles of bagging and ensembling can, in theory, be integrated into existing LLM-driven frameworks. Nevertheless, implementing such modifications would double the number of required inferences per sample, sharply escalating computational demand and, thereby, reducing practical feasibility. 

\subsection{Analysis \& Discussion}
\label{sec:analysis}
\paragraph{Ablation study.} To determine the impact of core components on the overall performance, we conduct a series of ablation studies targeting: (1) \textsf{-Tuning}: excluding the weight adjustment in the linear model, with the logit derived as a simple sum of the binary features; (2) \textsf{-Ensemble}: removing the ensemble aggregation; (3) \textsf{-Description}: removing the inclusion of feature description; and (4) \textsf{-Reasoning}: excluding Step-1 of the reasoning instruction during rule generation.

Table~\ref{tab:ablation} provides the average change in AUC following each ablation, computed over all datasets. Across the board, the exclusion of any of the specified components consistently leads to diminished performance. Notably, as the shot count rises, the contribution of weight tuning becomes increasingly pronounced, implying that richer data improves rule weighting accuracy and highlights the importance of this tuning. In contrast, the effects of feature description and reasoning instructions are especially significant for lower shot scenarios, underscoring the value of effectively leveraging the LLM’s prior knowledge when data are scarce. Finally, the ensemble strategy introduced in our framework shows stable improvements across all experimental settings.

\vspace*{-0.12in}
\paragraph{Training \& inference time comparison.} We assess and contrast the training and inference durations across the evaluated approaches. Experiments utilize the Adult dataset with a 16-shot training set, and unless otherwise specified, a single A100 GPU is employed. The TabLLM method is run on four A100 GPUs to support model parallelism. Table~\ref{tab:cost} presents detailed runtime statistics. Importantly, \textsf{FeatLLM} achieves inference speeds that are relatively low and on par with standard tabular algorithms. Conversely, other approaches leveraging LLMs—including In-context learning, TABLET, and TabLLM—require substantially longer inference periods. For training, \textsf{FeatLLM} demonstrates similar efficiency to other LLM-based solutions, utilizing only 30 API calls; this minimal query load helps ensure cost-effectiveness and can also be distributed for greater parallelism.

\vspace*{-0.12in}
\paragraph{Quality of features generated by \textsf{FeatLLM}.}
A straightforward experiment evaluates the utility of rule-based features produced by \textsf{FeatLLM} for real-world tasks. We compute the mean AUROC between these derived features and the true labels, and determine the proportion of features that attain an AUROC exceeding 0.5 on each dataset. 
Results are summarized in Table~\ref{tab:quality}, which indicates that the bulk of the generated rules have a significant association with the target variable and deliver valuable signals for solving the target problem (see “Before tuning” column). Furthermore, when training a linear model, rules with marginal relevance (identified by learned weights falling below a set threshold) are systematically excluded (see “After tuning” column). The threshold for filtering was established at 0.1 based on the distribution of weight values.

\vspace*{-0.12in}
\paragraph{Effect of adding spurious correlations.} To evaluate how well different approaches can mitigate the influence of irrelevant spurious correlations in few-shot scenarios, we carried out a focused analysis. Specifically, we used the Heart dataset and artificially appended randomly selected columns from the Adult dataset that are not pertinent to the predictive task. We assessed how model performance evolved as we incrementally introduced additional unrelated columns. For consistent evaluation, each method was provided with only 8 labeled samples, and we omitted methods such as SCARF or STUNT that rely on access to unlabeled data.

Figure~\ref{fig:spurious} presents how performance shifts with the introduction of spurious correlations. Among the methods, \textsf{FeatLLM} exhibited the smallest decline in accuracy attributable to irrelevant features. This resilience can be attributed to the framework’s emphasis on leveraging prior knowledge to explicitly connect features with the task during rule extraction. As a result, the model avoids relying on irrelevant columns, leading to more stable outcomes. Our observations indicate that rules involving noisy or irrelevant features made up merely 1-2\% of all extracted rules. Nevertheless, it is worth noting that \textsf{FeatLLM} may still inadvertently capture spurious associations and incorrectly infer causal relationships if supplementary columns originate from datasets with inherently related attributes (refer to Appendix~\ref{sec:spurious-appendix} for further discussion).

\vspace*{-0.12in}
\paragraph{Hyper-parameter analysis}
\textsf{FeatLLM} is primarily influenced by two hyper-parameters: the ensemble count and the number of rules per class retrieved in each LLM query. To assess their influence, we performed a series of experiments. Findings indicate that increasing the ensemble count enhances performance, although improvements plateau at approximately 20 ensembles (refer to Fig.~\ref{fig:appendix-hyperparameter1} in the Appendix). With respect to the rule count per class, although differences were not statistically significant, setting this value to 10 achieved the best balance between prompt length and predictive efficacy (as evidenced by Fig.~\ref{fig:hyperparameter2}). We hypothesize that while adding more rules can contribute richer information by increasing input dimensionality, it may also introduce overfitting, particularly in low-shot scenarios.

\section{Conclusion}
We introduce \textsf{FeatLLM}, a new approach that advances LLM-based few-shot tabular learning. Unlike conventional techniques that employ LLMs to directly predict outcomes for individual samples, \textsf{FeatLLM} leverages LLMs to generate predictive criteria, adopting a feature engineering perspective. By combining rule synthesis with a bagging-based ensemble mechanism, \textsf{FeatLLM} achieves efficient inference, operates exclusively via API calls without necessitating any training, and bypasses typical limitations related to feature dimensionality in data. As a result, \textsf{FeatLLM} delivers strong prediction accuracy, positioning it as an effective and data-efficient method for deploying LLMs on practical tabular datasets.

The current design of \textsf{FeatLLM} is optimized for low-shot learning cases. Future extensions will target feature creation in settings with larger data volumes, explore a broader range of feature types beyond rule-based features, and seek to enhance interpretability, thereby broadening the applicability to various domains.

\section{Broader Impact} The presented \textsf{FeatLLM} framework seeks to harness the latent knowledge within LLMs to elevate the effectiveness of tabular learning, surpassing the boundaries of traditional supervised methods. Our research contributes to improved data efficiency by leveraging LLM-based priors, which play a key role in reducing overfitting when annotated examples are scarce. \textsf{FeatLLM} is especially useful for professionals working in environments with substantial labeling barriers—such as those in the finance and healthcare sectors—where acquiring labeled data incurs significant cost or logistical difficulty, and the pretrained domain knowledge of LLMs can serve as a valuable supplement. Nevertheless, for widespread adoption, it is crucial for future investigations to examine how societal biases and potential misinformation contained within the LLM priors influence predictions, particularly since such priors could override or outweigh limited available labeled data.

\end{document}